
According to Sec. \ref{fpmatrixsection}, the electron Fokker-Planck equation
may be expressed in a general symbolic matrix form
\begin{equation}
\widehat{\Bbb{N}}X^{(k+1)}=Z^{(k)}  \label{generalmatrixform}
\end{equation}
where 
\begin{eqnarray}
\widehat{\Bbb{N}} &=&\frac{\widehat{\Bbb{A}}}{\Delta t}+\widehat{\Bbb{B}} \\
Z^{(k)} &=&\left( \frac{\widehat{\Bbb{A}}}{\Delta t}\right) X^{(k)}+\widehat{%
C}\left( X_{M},X^{(k)}\right) 
\end{eqnarray}
for the fully implicit time difference scheme, and 
\begin{eqnarray}
\widehat{\Bbb{N}} &=&\frac{\widehat{\Bbb{A}}}{\Delta t}+\frac{\widehat{\Bbb{B%
}}}{2} \\
Y^{(k)} &=&\left( \frac{\widehat{\Bbb{A}}}{\Delta t}-\frac{\widehat{\Bbb{B}}%
}{2}\right) X^{(k)}+\widehat{C}\left( X_{M},X^{(k)}\right) 
\end{eqnarray}
for the the Crank-Nicholson time scheme. 

The formal expression (\ref{generalmatrixform}) may be also used for the
electron drift kinetic equation as shown in Sec. \ref{dkematrixsection}, and
in this case, 
\begin{eqnarray}
\widehat{\Bbb{N}} &=&\widehat{\Bbb{G}} \\
Z^{(k)} &=&\widehat{\Bbb{G}}Y^{(k)}+\widehat{C}\left( X_{M},Y^{(k)}\right) +%
\widehat{\widetilde{W}}
\end{eqnarray}

In all cases, the matrices are square with similar structures, i.e. non-zero
elements are aligned along several diagonals which are roughly symmetricaly
placed around the main one\footnote{%
The lack of symmetry between upper and lower diagonals arises usually from
boundary conditions.}. 

The determination of $X^{(k+1)}$ requires to invert the system of equation (%
\ref{generalmatrixform}). The usual method based on a direct inversion by
the well known Gaussian elimination is ruled out, since the number of
operations for each direction inversion by Gaussian elimination is $\mathcal{%
O}\left( N^{3}\right) ,$ where $N$ is the number of rows (or columns) of
matrix $\widehat{\Bbb{N}}.$ Indeed, with $N$ that may reach $800000,$ the
number of operations become rapidly prohibitive. In addition, memory
requirements grow as $\mathcal{O}\left( N^{2}\right) $ since matrix $%
\widehat{\Bbb{N}}^{-1}$ is not sparse as matrix $\widehat{\Bbb{N}}$, leading
also to unacceptable storage constraints. 

Consequently, alternative methods must be used, in order to perform a fast
matrix inversion with a reduced storage capacity. The basic idea is to
factorize $\widehat{\Bbb{N}}$ into a upper and lower triangular matrices $%
\widehat{\Bbb{U}}$ and $\widehat{\Bbb{L}}$ respectively 
\begin{equation}
\widehat{\Bbb{N}}=\widehat{\Bbb{L}}\widehat{\Bbb{U}}
\end{equation}
as shown in Fig. \ref{figfactmatrix}, that are themselves highly sparse
matrices. A strong reduction of the computational effort may then be
foreseen, since the number of operations is considerably lower. This elegant
approach has been first considered for Fokker-Planck calculations for a five
diagonals operator as shown in Ref. \cite{sho94}, since in that case $%
\widehat{\Bbb{U}}$ and $\widehat{\Bbb{L}}$ are both triangular and
tridiagonal matrices. However, as discussed in Ref. \cite{pey98}, the fact
that cross-derivatives are considered explicitely with respect to the time
differencing scheme which puts strong limitations on the integration time
step $\Delta t,$ since it must be much lower than one, for an accurate
determination of the solution of the set of linear equations. Consequently,
the overall time duration for calculating the steady-state solution is very
long, an important drawback when the kinetic solver must be incorporated in
a chain of codes for realistic tokamak simulations. 

For this purpose, this method has be succesfully extended to the nine
diagonals operator case, allowing the use of much larger $\Delta t$ values
while the numerical scheme remains still stable, as shown in Ref. \cite
{pey98}. However, this approach is only useful for local kinetic
calculations, where trapped particle contribution is negligible, i.e. close
to the plasma center. When off-axis electron dynamics must be described, one
must in that case fully consider both circulating and trapped electron
dynamics which leads to a number of diagonals with non-zero elements that is
much larger, fifteen diagonals at each radial location as shown in Fig. \ref
{fpmatrix}. Therefore, the method developped for nine diagonal operators may
no more be used, since the number of operations to determine $\widehat{\Bbb{U%
}}$ and $\widehat{\Bbb{L}}$ increases dramaticaly. The fully 3-D approach
with radial transport makes this method unusable. 

However, an equivalent form of the approximate matrix factorization may
still be employed, noting that exact $\widehat{\Bbb{U}}$ and $\widehat{\Bbb{L%
}}$ matrices are made of a large number of very small coefficients, close to
zero. This indicates that information on the electron dynamics is fully
contained in a few set of coefficients, whose number is several order of
magnitude lower than the total one. Such a structure results in fact from
the well conditionning of the matrix $\widehat{\Bbb{N}},$ a physical
consequence that collisions predominate over all other physical processes.
Consequently, the general approach is to introduce a drop tolerance
parameter $\delta _{lu}$, below which small coefficients of exact $\widehat{%
\Bbb{U}}$ and $\widehat{\Bbb{L}}$ matrices are forced to zero. Since
coefficients of $\widehat{\Bbb{N}}$ lie between $0$ and $1$ because of the
main diagonal preconditioning, 
\begin{equation}
0\leq \delta _{lu}\leq 1
\end{equation}

When $\delta _{lu}$ is chosen close to $1$, approximate and exact matrices $%
\widehat{\Bbb{U}}$ and $\widehat{\Bbb{L}}$ are roughly equivalent, while
they differ strongly when $\delta _{lu}$ tends to zero. 
